\section{Design Science Research Methodology}
\label{sec:dsr_methodology}

Unlike traditional natural science methodologies that aim to discover and explain existing phenomena, Design Science Research (DSR) is rooted in pragmatic epistemology \cite{simon_1969}. DSR focuses on creating and evaluating innovative artifacts that solve identified operational problems, extending human and organizational capabilities while contributing prescriptive knowledge to the scientific body of knowledge \cite{hevner_2004, gregor_2013}.

In this research, we adopt the formal Design Science Research Methodology (DSRM) process model articulated by Peffers et al. \cite{peffers_2007}. As summarized in Table~\ref{tab:dsr_elements}, our research cycle operationalizes the six sequential DSRM activities:

\begin{table*}[!htbp]
\centering
\small
\caption{Design Science Research (DSR) Framework Mapping for the PANOPTES Reference Architecture}
\label{tab:dsr_elements}
\begin{tabularx}{\textwidth}{@{}l p{4.2cm} X@{}}
\toprule
\textbf{DSR Dimension} & \textbf{Epistemological Definition} & \textbf{Operationalization in PANOPTES} \\ \midrule
\textbf{Problem Domain} & Real-world operational dilemma in cloud-native microservices & Elevated Mean Time to Diagnose ($MTTDiag$), diagnostic inaccuracy under cascading faults, and unquantified probe overhead in distributed container environments (Wicked Problem). \\ \addlinespace[2pt]
\textbf{Class of Problems} & Generalized software engineering challenge & Automated, real-time correlation across heterogeneous, asynchronous telemetry modalities (metrics, logs, traces) without degrading host/cgroup performance. \\ \addlinespace[2pt]
\textbf{Artifact Type} & Prescriptive scientific contribution (Vaishnavi \& Kuechler) & \textbf{Model / Architecture:} Structural 4-layer reference framework with formal semantic bindings; \textbf{Method:} Span self-time ($T_{\text{self}}$) decomposition and exemplar correlation; \textbf{Instantiation:} Kubernetes IaC stack (Prometheus, Loki, Tempo, eBPF, OTel Collector). \\ \addlinespace[2pt]
\textbf{Design Principles} & Foundational engineering heuristics (Hevner et al., Walls et al.) & \textbf{Control Principles (\texttt{CP1}--\texttt{CP2}):} Non-invasive causal linkage ($\le 10\%$ CPU) and unified semantic representation (OTLP/Exemplars); \textbf{Contingency Heuristics (\texttt{ConH1}--\texttt{ConH3}):} Declarative control (IaC/CRDs), configuration drift rejection, and adaptive feedback to contain the Observer Effect. \\ \addlinespace[2pt]
\textbf{Evaluation Rigor} & Multi-criteria artifact evaluation (Hevner et al.) & Technical feasibility and overhead benchmarking (CP1 verification) coupled with architectural utility validation across $N=30$ stratified failure scenarios on CNCF Astronomy Shop. \\ \bottomrule
\end{tabularx}
\end{table*}


\begin{enumerate}[leftmargin=1.2em, itemsep=3pt, topsep=4pt]
    \item \textbf{Problem Identification and Motivation:} As characterized in Section~\ref{sec:background}, diagnosing distributed cloud-native failures is hindered by telemetry silos, the Cascading Span Bias, and uncontrolled probe overhead. These challenges constitute a ``wicked problem'' characterized by complex interdependencies and conflicting trade-offs between diagnostic accuracy and resource overhead.
    \item \textbf{Define Objectives of a Solution:} The objectives are quantitatively formalized through our three Research Questions (RQ1--RQ3): (i) establish an architecturally unified, 4-layer reference framework with formal semantic invariants; (ii) achieve $>80\%$ automated root-cause accuracy under multi-class chaos perturbations; and (iii) constrain total observability CPU overhead to $\le 10\%$ of cluster capacity.
    \item \textbf{Design and Development (The Artifact):} Following the artifact taxonomy of Vaishnavi and Kuechler \cite{vaishnavi_2015}, PANOPTES delivers three complementary contributions:
    \begin{itemize}[leftmargin=1.0em, itemsep=1pt]
        \item \textit{Model / Reference Architecture:} A 4-layer conceptual model establishing clean separation between ingestion, storage, correlation, and remediation.
        \item \textit{Method / Algorithms:} A deterministic causal correlation algorithm based on exclusive Span Self-Time decomposition ($T_{\text{self}}$) and multimodal evidence fusion.
        \item \textit{Instantiation:} A fully automated Kubernetes-native deployment pipeline codified in Ansible playbooks, Helm charts, and custom OpenTelemetry Collector pipelines.
    \end{itemize}
    \item \textbf{Demonstration:} We instantiate the PANOPTES reference architecture within an isolated, multi-node Kubernetes testbed running the CNCF OpenTelemetry Astronomy Shop benchmark, simulating real-world e-commerce microservice traffic under Locust user load.
    \item \textbf{Evaluation (Artifact Validation \& Benchmarking):} Following the evaluation taxonomy of Hevner et al. \cite{hevner_2004} and Gregor and Hevner \cite{gregor_2013}, we evaluate the artifact along three complementary dimensions: (i) \textit{Technical Feasibility and Invasiveness} by benchmarking the steady-state CPU and memory overhead against Control Principle \textbf{CP1}; (ii) \textit{Methodological Soundness and Diagnostic Utility} by comparing the PANOPTES correlation method ($T_{\text{self}}$ and multimodal evidence fusion) against unimodal baselines across $N=30$ stratified validation campaigns; and (iii) \textit{Architectural Completeness} by verifying closed-loop telemetry adaptation under resource pressure.
    \item \textbf{Communication:} The scientific contributions, reference models, and empirical findings are documented in this paper and released as an open-access replication package.
\end{enumerate}

To ensure architectural transparency and traceability throughout the design cycle, all critical technical design choices and trade-offs were codified using Markdown Architectural Decision Records (MADRs), adhering to the standard template proposed by Kopp et al.
